\documentclass[11pt]{article}
\usepackage[a4paper,margin=25mm]{geometry}
\usepackage[T1]{fontenc}
\usepackage{lmodern}
\usepackage{booktabs}
\usepackage{amsmath}
\usepackage{graphicx}
\usepackage{hyperref}
\usepackage{xcolor}
\usepackage{listings}
\lstset{basicstyle=\ttfamily\small, columns=fullflexible, keepspaces=true,
  breaklines=true, frame=single, framerule=0.3pt, rulecolor=\color{gray}}
\newcommand{\code}[1]{\texttt{#1}}
\newcommand{\us}{\textmu s}

\title{\code{wavefit.py}: the slow wave per channel and a baseline that follows it}
\author{BGOTSCAN}
\date{2026-10-15}

\begin{document}
\maketitle

\section{Purpose}

\code{cmnoise.py} found the noise of the charge to be a slow wave near
0.33~MHz, a 3~\us{} period, present in every channel, dead or live, with
its own phase per channel, and in the live channels a slower drift on
top (\code{cmnoise\_doc.tex}). The charge of a channel over the 5~\us{}
window, with the baseline from the last 2~\us{} of the record, carries
2400~ADC\,ns of this against 700 for white noise. A running baseline
cannot follow a 3~\us{} wave without following the 2.7~\us{} pulse, but
a baseline model fitted on the pulse-free samples around the window and
evaluated inside it can, if the wave keeps its phase over the window.
This script measures the wave where there is no signal, the tail of the
record in every channel and the whole record in the dead channels, and
tests the baseline model on an empty window: the scatter of the charge of
a window with nothing in it, with the current baseline and with the
fitted one, is the gain of the method. Nothing here touches a pulse.

\section{Usage}

\begin{lstlisting}
python3 src/py/wavefit.py RUN [--events 0-299] [--tail 12] [--window 1,5] [--before 1] [--smooth 62]
        [--band 0.2,1.0] [--dead LIST] [--offline FILE] [--examples 3] [--save NAME]
\end{lstlisting}

\begin{center}
\begin{tabular}{l l p{78mm}}
\toprule
option & default & meaning \\
\midrule
\code{RUN} & & run name or path of the raw file \\
\code{--events a-b} & \code{0-299} & events read (0-based) \\
\code{--tail T} & 12 & the tail: the last $T$~\us{} of the record \\
\code{--window S,L} & \code{1,5} & the empty window: starts $S$~\us{} into the tail, lasts $L$~\us \\
\code{--before B} & 1 & samples before the window used by the fit, in \us{} (the pre-trigger of a real pulse) \\
\code{--after A} & 2 & samples right after the window averaged for the ``after'' baseline, in \us \\
\code{--smooth N} & 62 & running-mean width in samples (0.5~\us) for the autocorrelation \\
\code{--band a,b} & \code{0.2,1.0} & MHz band where the wave's frequency is searched \\
\code{--dead LIST} & the twelve of \code{chanstat\_doc.tex} & dead channels \\
\code{--offline FILE} & \code{pass/RUN\_offline.h5} & per-event classes; without it every event is used \\
\code{--examples N} & 3 & events drawn in the by-eye figure \\
\code{--save NAME} & & save the figures as \code{src/py/plots/RUN/NAME\_wavefit\_*.png} \\
\bottomrule
\end{tabular}
\end{center}

\section{Definitions}

Waveforms minus the WebDump baseline of the channel; the events are the
REAL events without pile-up of \code{offline.py} whose maximum lies more
than 4~\us{} before the tail, as in \code{cmnoise.py}. The tail is the last
$T$~\us{} of the record. Inside it, the empty window $W$ runs from $S$ to
$S + L$~\us{} after the start of the tail; the fit region $F$ is the $B$~\us{}
before the window and everything after it to the end of the record; the
end-of-record baseline region $E$ is the last 250 samples, as in the pass.

\paragraph{Autocorrelation.} Per channel and event, the tail smoothed
with a running mean of $N$ samples, minus its mean; its autocorrelation
$\rho(\tau)$ for lags $\tau$ from 0 to 10~\us, normalised to 1 at zero
lag; averaged over events and over the channels of each class, dead and
live. The white noise contributes a spike of the width of the smoothing
at zero lag; the wave gives lobes alternating in sign every half period,
as many as the periods over which it keeps its phase; a drift gives one
lobe. Figure \code{autocorr}: left, $\rho(\tau)$ for the two classes;
right, the rms of the running mean of one channel against its width,
from one sample to 5~\us, averaged over channels and 60 events, with the
white-noise line $\sigma/\sqrt N$ (the counterpart of the toy figure of
\code{noisetoy\_doc.tex}). The terminal gives, per class, the lag of the first zero
crossing, the depth of the first negative lobe and the height of the
second positive lobe, the quantities that say how far the tail can
predict the wave.

\paragraph{Frequency.} Per channel, the power spectrum of the raw tail
averaged over events, and the frequency of its maximum within the band
$[a, b]$~MHz, above the low-frequency rise and below the white floor.
Figure \code{spectra}: all channels overlaid, dead in red, live in blue,
and the $8\times8$ map of the peak frequency. The run's wave frequency
$f_0$ is the median over the live channels and is the frequency used by
the baseline model.

\paragraph{Baselines of the empty window.} Four charges of $W$ per
channel and event, all in ADC\,ns:
\begin{align*}
  q_E &= \Delta t \sum_{i \in W} (x_i - \bar x_E), \\
  q_A &= \Delta t \sum_{i \in W} (x_i - \bar x_A), \\
  q_{\rm line} &= \Delta t \sum_{i \in W} (x_i - a - b\,t_i), \\
  q_{\rm wave} &= \Delta t \sum_{i \in W} \big(x_i - a - b\,t_i - c \sin 2\pi f_0 t_i - d \cos 2\pi f_0 t_i\big),
\end{align*}
with $\bar x_E$ the mean over $E$ (the pass), $\bar x_A$ the mean over
the $A$~\us{} right after the window, and $a, b$ or $a, b, c, d$
from a linear least-squares fit of the model to the samples of $F$
(the window itself excluded). The scatter of each charge over the
events, $\sigma_E$, $\sigma_{\rm line}$, $\sigma_{\rm wave}$, is the
noise that method would leave on a real charge, since the window is
empty; the gain of a method is $\sigma_E / \sigma$. The white-noise
floor is $\sigma_{\rm w} \sqrt{n_W}\,\Delta t$ with $\sigma_{\rm w}$ the
clipped rms of the channel's samples, 700~ADC\,ns for 4 counts and 625
samples, raised by the noise of the baseline itself.

\paragraph{Figures.} \code{examples}: for \code{--examples} events, one
dead and one live channel, the tail with the window shaded, the
end-of-record level as a flat line and the two fitted models drawn
through the window, so that one sees whether a model follows the wave
into the window. \code{hist}: the histograms of $q_{E+\rm cm}$, $q_E$, $q_{\rm line}$,
$q_{\rm wave}$ over channels and events, dead and live separately, with
their scatters. \code{gainmap}: the $8\times8$ map of
$\sigma_E / \sigma_{\rm wave}$ per channel. \code{sum}: the same four
charges for the sum of the live channels (the per-channel residuals
summed), where the weak common part of the noise and the cost of the
dead-channel common mode show.

\paragraph{The pass's common mode.} A fourth charge, $q_{E+\rm cm}$: the
125-sample running mean of the dead-channel mean subtracted from every
channel, as the pass does with \code{--cm dead}, then the end-of-record
baseline. Its scatter on the sum is compared with the noise of the dead
channels' mean times the number of live channels, $n_l\,\sigma_E^{\rm
dead}/\sqrt{n_d}$, added in quadrature to the raw sum.

\paragraph{Terminal.} Per class: the four scatters, the white floor and
the gains, for the channels and for the sum; $f_0$ and its spread
over channels; the autocorrelation numbers above.

\section{Findings (2026-10-15, 195K\_41V\_th20 and 197K\_39V\_th20, 300 events each)}

\begin{itemize}
\item \emph{There is no wave.} The single-channel spectra are a smooth
  rise below 1~MHz with no line in any channel, dead or live; the peak
  search in the band returns its lower edge (0.25~MHz, the third
  frequency bin). The autocorrelation has one lobe, zero at 1.5~\us, and
  nothing after: the slow noise forgets itself within 2~\us. The
  ``0.33~MHz oscillation'' read off the mean of the channels in
  \code{cmnoise\_doc.tex} was a feature of a noisy spectrum of the mean
  at 0.08~MHz resolution and of the eye; withdrawn. The slow noise is
  broadband, $1/f$-like, per channel.
\item \emph{Dead and live channels carry the same noise.} Empty-window
  charge with the end-of-record baseline: 3473~ADC\,ns for the dead
  channels, 3581 for the live, against a white floor of 756, at both
  biases. The statement of \code{chanstat\_doc.tex} that the dead channels
  show the charge noise of white noise is contradicted by this direct
  measurement and should be rechecked there.
\item \emph{A baseline fitted around the window gains a factor 1.37}, a
  line through the 1~\us{} before and the 6~\us{} after: 3500 to 2550 per
  channel; the mean of the 2~\us{} right after the window gains 1.14
  (3050); a line fitted on the after side alone extrapolates its slope
  into the window and loses (5100). The sinusoid adds nothing (2600 to
  2770), as a broadband noise cannot be predicted inside the window
  from outside. The residual is 3.4 times the white floor: the $1/f$
  noise inside the window itself.
\item \emph{The dead-channel common mode of the pass is the noise of
  the sum.} The raw sum of the 52 live channels over the empty window
  scatters by 29\,000 at 41~V and 36\,000 at 39~V, against 25\,800 for
  independent channels: the common part is small. After subtracting the
  running mean of the 12 dead channels, as the pass does, the scatter is
  60\,000 and 65\,000: the dead channels' own slow noise, $3473/\sqrt{12}
  \times 52 = 52\,000$, added in quadrature (59\,700 and 63\,600
  predicted). This is the 66\,000 of \code{offline.py}. The dead channels
  share nothing with the live ones within the record (correlation 0.00),
  so the subtraction removes nothing and adds their noise; per channel it
  costs 4\% (3581 to 3719), on the sum a factor 2.
\end{itemize}

Consequence: drop the common-mode subtraction (\code{--cm none} in the
pass) and take the per-channel baseline from a line fitted around the
window. The sum's noise then falls from 66\,000 to about 27\,000, and the
per-channel noise from 3700 to 2600. The $1/f$ noise within the window
remains, 3.4 times white, and is a hardware question for the front-end.

\section{Examples}

\begin{lstlisting}
python3 src/py/wavefit.py 195K_41V_th20 --save wf195
python3 src/py/wavefit.py 197K_39V_th20 --save wf197
\end{lstlisting}

\end{document}
